\subsection{Numerical evaluation and comparison protocol}
\label{subsec:numerical_evaluation_protocol}

The numerical assessment is conducted on three parametrized incompressible-flow
problems exhibiting steady, Hopf-onset, and periodic dynamics: flow past a
circular cylinder, flow past a centered square cylinder, and the fluidic
pinball. These examples involve different geometries, discretizations, and
regime-transition ranges, and are used collectively to assess the predictive
performance and robustness of the proposed RAL-MoE-ROM framework.

The numerical experiments are organized to examine two main components of the
method. First, the regime-local specialists are evaluated in their respective
parameter domains to assess the accuracy and stability of the local reduced
dynamics over autonomous multi-step rollouts. Where compatible baselines are
available, the proposed specialists are compared with alternative neural
architectures, data-driven reduced models, classical POD-Galerkin models and
global reduced representations. These comparisons and ablation studies are
used to assess the contributions of regime localization, the projected
Galerkin dynamics, and the learned closure. Second, the hierarchical assembly
is evaluated in the overlapping parameter regions. Individual adjacent
specialists and hard routing are compared with fixed convex combinations,
router-probability blending, the proposed T2-C correction, and, where
available, a per-window convex oracle. These experiments quantify whether
adjacent local models provide complementary predictions and whether the
proposed fusion mechanism can exploit this complementarity without modifying
the underlying specialist dynamics.

The three benchmark problems play complementary roles in this assessment. The
flow pass cylinder example examines the regime-local specialists and their
long-horizon dynamical behaviour across the steady-Hopf-periodic transition.
The flow pass a centered square example evaluates the multi-chart construction
for a different geometry and investigates the steady-hopf and hopf-periodic fusion mechanisms.
The fluidic-pinball example further compares the proposed local learned
dynamics with a classical POD-Galerkin/Pressure-Poisson ROM and a global
data-driven reduced model POD-Deeponet. The baseline sets are therefore not required to be
identical across the three examples; rather, each comparison is used only for
the methodological component for which a consistent numerical evaluation is
available.

All training, validation, and held-out test sets are partitioned by complete
parameter trajectories. Model selection and all router and fusion-model
selection are performed using training and validation data only. Unless stated
otherwise, $K$ denotes the number of successive prediction steps in the
corresponding snapshot sequence. Since the snapshot interval may differ among
the numerical examples, the associated physical-time horizon is reported
separately when required.

We use a common set of weighted relative errors to evaluate the reconstructed
physical fields. Let $\boldsymbol{u}_n$ and $p_n^\circ=\Gamma_p(p_n)$ denote
the reference velocity and gauge-corrected pressure at an evaluation time, and
let $\widehat{\boldsymbol{u}}_n$ and $\widehat{p}_n^\circ$ denote the
corresponding predictions. Using the discrete physical inner products defined
by the quadrature matrices $M_u$ and $M_p$, the relative velocity and pressure
errors are
\begin{align}
E_u
&=
100\times
\frac{
\left\|
\widehat{\boldsymbol{u}}_n-\boldsymbol{u}_n
\right\|_{M_u}
}{
\left\|
\boldsymbol{u}_n
\right\|_{M_u}+\varepsilon
},
&
E_p
&=
100\times
\frac{
\left\|
\widehat{p}_n^\circ-p_n^\circ
\right\|_{M_p}
}{
\left\|
p_n^\circ
\right\|_{M_p}+\varepsilon
}.
\label{eq:common_field_errors}
\end{align}

The velocity and pressure errors are always retained separately in addition to
$E_{\mathrm{joint}}$.

For predictions evaluated within the same regime-local POD chart
$r\in\{S,H,P\}$, we additionally measure errors in the reduced coordinates,
\begin{align}
E_a^{(r)}
&=
100\,
\frac{
\left\|
\widehat{\boldsymbol{a}}^{(r)}
-\boldsymbol{a}^{(r)}
\right\|_2
}{
\left\|
\boldsymbol{a}^{(r)}
\right\|_2+\varepsilon
},
&
E_b^{(r)}
&=
100\,
\frac{
\left\|
\widehat{\boldsymbol{b}}^{(r)}
-\boldsymbol{b}^{(r)}
\right\|_2
}{
\left\|
\boldsymbol{b}^{(r)}
\right\|_2+\varepsilon
}.
\label{eq:common_coordinate_errors}
\end{align}
These quantities characterize errors in the local fluctuation dynamics and are
used only within a fixed POD coordinate system; reduced-coordinate errors from
different local charts are not compared directly. Since the physical-field
errors include the regime-local mean whereas the POD coefficients describe
mean-subtracted fluctuations, a small value of $E_u$ or $E_p$ does not
necessarily imply a comparably small value of $E_a^{(r)}$ or $E_b^{(r)}$.

In addition to these finite-horizon errors, regime-dependent dynamical
diagnostics are used where appropriate to assess autonomous stability and
long-time behaviour, including boundedness, fixed-point behaviour, oscillation
amplitude and frequency, phase accuracy, and orbit geometry. These diagnostics
complement the field and reduced-coordinate errors, since accurate one-step or
finite-horizon prediction alone does not guarantee faithful autonomous
long-time dynamics.




\subsection{Flow past a circular cylinder}
\label{subsec:flow_past_circular_cylinder}

The flow past a circular cylinder is used to evaluate RAL-MoE-ROM across
steady, Hopf-onset, and periodic flow regimes.  The numerical study examines
both the long-horizon predictive accuracy of the regime-local reduced models and the effectiveness of the hierarchical routing and fusion mechanism in the overlapping parameter regions.

The high-fidelity model solves the two-dimensional incompressible
Navier--Stokes equations for a cylinder of diameter $D=1$ with kinematic
viscosity $\nu=10^{-3}$.  The inflow speed is
$U_\infty=Re\,\nu/D$.  The computational domain is
$[-10,20]\times[-10,10]$, with the cylinder centred at the origin and having
radius $0.5$; a single cell layer over
$z\in[-0.05,0.05]$ realizes the two-dimensional calculation.  A prescribed
velocity is imposed at the inlet, a zero-normal-gradient condition at the
outlet, no slip on the cylinder, and slip on the upper and lower boundaries.
Pressure is fixed to zero at the outlet and uses zero-normal-gradient
conditions elsewhere.  All three reduced databases share the same ordered
grid of 97,368 points and the same lumped point-area weights.  Pressure is
placed in a common gauge by subtracting its area-weighted spatial mean from
each snapshot before projection, reconstruction and evaluation.
\begin{table*}[htbp]
\centering
\caption{Data partitions and rollout horizons.
The numbers of training, validation, and held-out test samples refer to
complete Reynolds-number trajectories. The physical-time horizon is reported
as the range associated with the corresponding $K$-step rollouts.}
\label{tab:cylinder_data_contracts}

\scriptsize
\begin{adjustbox}{max width=\textwidth}
\begin{tabular}{@{}p{0.11\textwidth}p{0.17\textwidth}c
                p{0.16\textwidth}p{0.14\textwidth}p{0.18\textwidth}@{}}
\toprule
Regime
& Reynolds-number range
& Train/validation/test trajectories
& Training snapshots
& Rollout horizon
& Physical-time horizon \\
\midrule

Steady
& $20.0 \le Re \le 46.4$
& $14/2/4$
& 894
& $K=56$
& $234.604$--$431.274$ \\

Hopf
& $45.5 \le Re \le 59.2$
& $29/2/3$
& 1935
& $K=56$
& $1179.953$--$1309.781$ \\

Periodic
& $57.3 \le Re \le 200$
& $53/6/4$
& 8839
& $K=48$
& $177.959$--$598.701$ \\

\bottomrule
\end{tabular}
\end{adjustbox}
\end{table*}
The Steady, Hopf, and Periodic datasets cover
$20.0\le Re\le46.4$, $45.5\le Re\le59.2$, and
$57.3\le Re\le200$, respectively.
Table~\ref{tab:cylinder_data_contracts} summarizes the corresponding
training, validation, and held-out test partitions, all defined at the level
of complete parameter trajectories. For each regime, the local velocity and pressure POD spaces are constructed
with $r_u=r_p=32$. The reported rollout length $K$ denotes the number of successive steps in the
corresponding native snapshot sequence. Because the sampling intervals differ
among the three regime-specific databases, equal values of $K$ do not in
general correspond to equal physical-time horizons; the associated time spans
are therefore reported separately in
Table~\ref{tab:cylinder_data_contracts}.


The proposed regime-local specialists are compared with three ablation
models targeting different components of RAL-MoE-ROM. Vanilla-FNN-MoE
replaces the structured nonlinear-affine bilinear expert parameterization
with standard feed-forward networks while retaining the remaining
regime-local modeling components. DataOnly-MoE removes the explicit
POD-Galerkin velocity dynamics and Pressure-Poisson pressure reconstruction
and instead learns the reduced-state update directly from the temporal
history. These two ablations assess, respectively, the contributions of the
structured expert architecture and the physics-based reduced dynamics. Global MoE is used to assess the effect of regime localization. It replaces
the three regime-specific POD spaces and specialists by a single global POD
representation and a single MoE model constructed from the pooled training
trajectories across the complete parameter range. 



The common field and reduced-coordinate errors are those in
Eqs.~\eqref{eq:common_field_errors}--\eqref{eq:common_coordinate_errors}.
The cylinder case additionally evaluates pressure fixed-point error,
pressure drift and paired-perturbation gain in the Steady regime, and
amplitude, frequency, phase, normalized orbit distance, velocity and pressure
energy drift, finite fraction, divergent windows and false growth in the
Hopf and Periodic regimes.  The complete threshold definitions are given in
the supplementary material.  These diagnostics test asymptotic behaviour
separately from finite-horizon field and coordinate accuracy.

% Suggested placement: immediately after the setup, baseline and metric
% paragraphs.  The file defines Tables 1--4 proposed for this subsection.
\input{drafts/cylinder_wake_20260726/tables_cylinder_wake}
\begin{table*}[htbp]
\centering
\caption{Held-out long-horizon prediction errors for the circular-cylinder
case. Each interval reports the minimum--maximum value over the held-out
Reynolds-number trajectories of the corresponding regime. N/E indicates that
held-out evaluation was not performed because the validation-stage stability
requirement was not satisfied.}
\label{tab:cylinder_specialist_long}

\scriptsize
\setlength{\tabcolsep}{5.0pt}
\begin{adjustbox}{max width=\textwidth}
\begin{tabular}{@{}llcccc@{}}
\toprule
Regime
& Model
& $E_u$ (\%)
& $E_p$ (\%)
& $E_a$ (\%)
& $E_b$ (\%) \\
\midrule

Steady
& Proposed specialist
& 0.0330--0.1723
& 0.2500--1.3249
& 0.2480--0.3858
& 6.0857--9.4108 \\

& Vanilla-FNN-MoE
& 0.0384--0.2995
& 2.2163--18.5603
& 0.2562--0.5373
& 6.0504--22.6792 \\

& DataOnly-MoE
& 0.1637--0.5122
& 4751.1025--27275.4645
& 0.8749--3.6893
& 15945.6192--58856.7261 \\

& Global MoE
& 0.0845--0.3404
& 1.5226--7.4363
& 20.7911--44.2564
& 25.2651--101.0296 \\

\midrule

Hopf
& Proposed specialist
& 0.0032--0.0288
& 0.0355--0.4321
& 0.0418--0.8479
& 0.2164--5.7742 \\

& Vanilla-FNN-MoE
& N/E
& N/E
& N/E
& N/E \\

& DataOnly-MoE
& 0.1260--0.3880
& 1280.2811--1524.5949
& 1.6428--18.0141
& 8810.2142--40976.6418 \\

& Global MoE
& 0.0110--0.0268
& 0.0560--0.2882
& 24.1395--94.7850
& 29.9005--78.2764 \\

\midrule

Periodic
& Proposed specialist
& 0.3809--1.3473
& 1.6825--5.9601
& 1.6843--2.9151
& 2.3345--6.7689 \\

& Vanilla-FNN-MoE
& 1.0393--3.0836
& 3.9659--13.3257
& 4.4136--9.4908
& 5.3593--19.0759 \\

& DataOnly-MoE
& 0.7854--1.8677
& 2.2511--10.8189
& 2.2001--4.1827
& 3.4416--5.3181 \\

& Global MoE
& 0.5873--1.1726
& 2.1896--6.3566
& 5.3269--16.9732
& 6.6735--21.9848 \\

\bottomrule
\end{tabular}
\end{adjustbox}
\end{table*}
Table~\ref{tab:cylinder_specialist_long} compares the held-out
long-horizon prediction errors of the proposed regime-local specialists and
the three ablation models. The results reveal a marked dependence of the
ablation effects on the underlying flow regime. For Global MoE,
$E_a$ and $E_b$ are evaluated in its own global POD coordinates and are
therefore used only as diagnostics of its reduced-state prediction; they are
not directly compared with the coefficient errors associated with the
regime-local POD spaces.

In the Steady regime, the proposed specialist gives the lowest
physical-field error range among the considered models. Its velocity error is
$0.0330$--$0.1723\%$, compared with $0.0384$--$0.2995\%$ for
Vanilla-FNN-MoE and $0.0845$--$0.3404\%$ for Global MoE. A more pronounced
difference is observed in pressure, where the proposed specialist yields
$E_p=0.2500$--$1.3249\%$, whereas the upper bounds increase to
$18.5603\%$ for Vanilla-FNN-MoE and $7.4363\%$ for Global MoE.
Within the common regime-local representation, the proposed specialist also
maintains small velocity-coordinate errors and substantially reduces the
upper bound of the pressure-coordinate error relative to Vanilla-FNN-MoE.
By contrast, DataOnly-MoE produces pressure errors of
$4.75\times10^{3}$--$2.73\times10^{4}\%$, accompanied by extremely large
pressure-coordinate errors. These results indicate that the explicit
physics-based reduced dynamics are particularly important for pressure
prediction in the Steady regime. 


The differences are more pronounced in the Hopf regime. The proposed
specialist achieves velocity and pressure errors of
$0.0032$--$0.0288\%$ and $0.0355$--$0.4321\%$, respectively.
Vanilla-FNN-MoE is not reported on the held-out trajectories because its
selected model did not satisfy the validation-stage stability requirement.
Removing the physics-based reduced dynamics leads to a substantial loss of
accuracy: DataOnly-MoE increases the pressure error to
$1.28\times10^{3}$--$1.52\times10^{3}\%$ and the pressure-coordinate error
to approximately $8.8\times10^{3}$--$4.1\times10^{4}\%$. Global MoE, in contrast, retains small physical-field errors over the Hopf
test range, with $E_u=0.0110$--$0.0268\%$ and
$E_p=0.0560$--$0.2882\%$. Thus, the regime-local formulation does not
uniformly outperform the global model in the mean-containing physical-field
metrics for this regime. Nevertheless, the proposed specialist achieves
highly accurate prediction within its regime-specific reduced coordinates,
whereas the Global MoE coefficients belong to a different global POD space
and cannot be compared directly. The Hopf results therefore primarily
demonstrate the benefit of the hybrid Galerkin-closure formulation relative
to the architecture and data-only ablations, while showing that a global
train-only representation can remain competitive in physical-field accuracy
over this parameter range.



In the Periodic regime, the proposed specialist consistently improves upon
Vanilla-FNN-MoE in both the physical fields and the local reduced
coordinates. The velocity-error range decreases from
$1.0393$--$3.0836\%$ to $0.3809$--$1.3473\%$, while the pressure-error
range decreases from $3.9659$--$13.3257\%$ to
$1.6825$--$5.9601\%$. DataOnly-MoE is substantially more competitive here
than in the Steady and Hopf regimes: its velocity and pressure errors remain
within $0.7854$--$1.8677\%$ and $2.2511$--$10.8189\%$, respectively.
The proposed specialist nevertheless retains the lower overall physical-field
error ranges. Global MoE also provides competitive Periodic predictions. Its maximum
velocity error is slightly lower than that of the proposed specialist,
whereas the proposed specialist yields the lower pressure-error range.
Consequently, the advantage of regime localization is not uniform in the
finite-horizon field metrics. In contrast, the comparison with
Vanilla-FNN-MoE shows a clear benefit from the structured specialist
formulation, while the substantially improved DataOnly-MoE performance
relative to the other two regimes indicates that the importance of the
explicit physics-based closure is itself regime dependent.

Overall, the ablation study shows that the contributions of the individual
modeling components are regime dependent. The proposed specialist provides
a consistent improvement over the Vanilla-FNN-MoE architecture, with
particularly pronounced reductions in pressure error. Removing the explicit
Galerkin and Pressure-Poisson contributions severely degrades the Steady
and Hopf predictions, whereas the purely data-driven formulation remains
considerably more competitive in the Periodic regime. The Global MoE
comparison further shows that regime localization should not be interpreted
as providing uniform pointwise superiority over a global reduced
representation; rather, its benefit must be assessed together with the
regime-specific dynamical modeling and the subsequent hierarchical
combination examined below.



\begin{table}[htbp]
\centering
\caption{Held-out prediction errors in the $S$--$H$ overlap at $K=56$.
Each interval reports the minimum--maximum value over the three held-out
overlap trajectories.}
\label{tab:cylinder_sh_boundary}

\scriptsize
\setlength{\tabcolsep}{6pt}
\begin{tabular}{@{}lccc@{}}
\toprule
Model
& $E_u$ (\%)
& $E_p$ (\%)
& $\alpha_S$ \\
\midrule

E2 Top-1
& 0.0357--0.0410
& 1.6746--1.9028
& 1.0000 \\

T2-C
& 0.0148--0.0355
& 0.7387--1.7212
& 0.0861--0.9987 \\

Global MoE
& 0.0851--0.0953
& 1.9517--2.3106
& -- \\

\bottomrule
\end{tabular}
\end{table}
Table~\ref{tab:cylinder_sh_boundary} summarizes the held-out
$S$--$H$ overlap results at $K=56$. Relative to the hard E2 routing,
T2-C reduces the velocity-error range from
$0.0357$--$0.0410\%$ to $0.0148$--$0.0355\%$ and the
pressure-error range from $1.6746$--$1.9028\%$ to
$0.7387$--$1.7212\%$. Both errors are reduced at each of the three
held-out Reynolds numbers, indicating that the output-level combination
exploits complementary predictions of the adjacent Steady and Hopf
specialists over the tested overlap. The learned fusion does not collapse uniformly to either specialist.
The Steady weight varies from $\alpha_S=0.0861$ to $0.9987$ across the
held-out cases, ranging from an almost purely Steady prediction to a
predominantly Hopf contribution. Since the two specialists are advanced
independently in their own reduced coordinates and are combined only after
physical-field reconstruction, this improvement is obtained without
introducing a common latent state or feeding the fused solution back into
either local dynamical system. Global MoE is included as the train-only ablation of the regime-local
decomposition. Over the same held-out overlap parameters, its velocity and
pressure errors lie in the ranges $0.0851$--$0.0953\%$ and
$1.9517$--$2.3106\%$, respectively. The comparison therefore indicates
that, for this $S$--$H$ overlap, the regime-local construction combined
with T2-C provides lower physical-field errors than the single global
reduced model.

% Suggested placement: after the specialist and boundary quantitative
% discussion.  If space is limited, Fig. 2 can be moved to the supplement,
% while Figs. 1, 3 and 4 should remain in the main paper.


\begin{figure*}[htbp]
\centering

\includegraphics[width=0.85\textwidth]
{drafts/cylinder_wake_20260726/figures/periodic1.png}

\vspace{0.5em}

\includegraphics[width=0.85\textwidth]
{drafts/cylinder_wake_20260726/figures/periodic1_1.png}

\caption{Representative held-out Periodic prediction at
$Re=70.314636$ and $K=48$. The upper panel compares the CFD reference with
Vanilla-FNN-MoE, DataOnly-MoE, and the proposed Periodic specialist in
velocity magnitude and gauge-corrected pressure. The lower panel shows the
corresponding pointwise absolute errors
$\lvert\Delta\boldsymbol{u}\rvert$ and $\lvert\Delta p^\circ\rvert$.}
\label{fig:cylinder_periodic_fields_errors}
\end{figure*}




\begin{figure*}[htbp]
\centering

\includegraphics[width=0.85\textwidth]
{drafts/cylinder_wake_20260726/figures/SH_1.png}

\vspace{0.5em}

\includegraphics[width=0.85\textwidth]
{drafts/cylinder_wake_20260726/figures/SH_1_1.png}

\caption{Representative held-out prediction in the $S$--$H$ overlap at
$Re=43.9$ and $K=56$. The upper panel compares the CFD
reference, the two local specialists, and the fused prediction; the lower
panel shows the corresponding pointwise velocity and gauge-corrected pressure
errors. For this trajectory, T2-C assigns
$\alpha_S=0.1891$ and $\alpha_H=0.8109$, illustrating a nontrivial
combination of the two adjacent specialists.}
\label{fig:cylinder_sh_fields_errors}
\end{figure*}



The field visualizations in
Figs.~\ref{fig:cylinder_periodic_fields_errors} and
\ref{fig:cylinder_sh_fields_errors} complement the integral error measures
by resolving their spatial distribution. The Periodic example illustrates
the local-specialist comparison, whereas the $S$--$H$ example visualizes the
output-level combination of two independently evolved local reduced models.
These figures are used for spatial interpretation; the quantitative
comparisons remain those reported in
Tables~\ref{tab:cylinder_specialist_long} and
\ref{tab:cylinder_sh_boundary}.In the Periodic case, the proposed specialist retains the alternating
vortex-shedding pattern of the reference solution. The pointwise errors are
concentrated primarily in the near wake and around the transported coherent
structures, whereas the Vanilla-FNN-MoE discrepancy extends more strongly
along the vortex street.The $S$--$H$ visualization further shows how the two local predictions remain
independent before reconstruction and output-level fusion. The resulting
field remains spatially smooth across the overlap, with the local error
structure reflecting the relative contribution selected by the T2-C gate.







Overall, the circular-cylinder case validates both levels of the proposed
framework. The regime-local specialists provide accurate long-horizon
predictions, while the ablation study reveals regime-dependent contributions
from the structured expert architecture, the physics-based reduced dynamics,
and regime localization. In the $S$--$H$ overlap, T2-C further improves the
hard E2 routing through output-level combination of independently evolved
local models. The results therefore support the regime-local hierarchical
construction without implying uniform superiority of every modeling component
across all flow regimes.
